\metatitle{Schubert polynomial variations}
\metadescription{Variants and extensions of Schubert polynomials, including type B/C/D, skew, double, affine, involution, quantum, universal, and twisted Schubert polynomials.}
\metakeywords{Schubert polynomial, double Schubert polynomial, quantum Schubert polynomial, affine Schubert polynomial, type B Schubert polynomial}


\section[schubertVariations]{Schubert polynomial variations}

This page collects variants and extensions of Schubert polynomials.  The basic
type $A$ theory, divided-difference definition, pipe-dream models, multiplication
rules, and complexity results are covered on the main
\hyperref[schubert]{Schubert polynomial} page.


\subsection[schubertBCD]{Type B/C/D Schubert polynomials}

\begin{polydata}{schubertBCD}
  Name     & Type B/C/D Schubert polynomials \\
  Space    & All \\
  Basis    & True \\
  Rating   & 2 \\
  Bib      & FominKirillov1996 \\
  Year     & 1996 \\
  Symbol   & $\schubert^B_\omega(\xvec)$ \\
  Keywords & divided-difference  \\
  Category & Schubert \\
\end{polydata}


In type $B$, there are combinatorial analogues of Schubert polynomials
introduced in \cite{FominKirillov1996}.
They also introduce \hyperref[stanleySymBC]{type $B$ Stanley symmetric functions}.

Divided difference formulas for Schubert polynomials in types $B$, $C$ and $D$ were
introduced by \name{Sara Billey} and \name{Mark Haiman} \cite{BilleyHaiman1995}.
See also \name{Sara Billey}'s PhD thesis on Schubert polynomials for classical groups \cite{Billey1994Thesis}.

Type $B$, $C$ and $D$ double Schubert polynomials are defined in \cite{IkedaMihalceaNaruse2011}.
\name{Tomoo Matsumura} gives a tableau formula for vexillary Schubert
$P$-polynomials in type $C$ \cite{Matsumura2023}.  The formula uses flagged
factorial Schur $Q$-functions and a tableau model adapted to vexillary signed
permutations.
\name{Evgeny Smirnov} and \name{Anna Tutubalina} construct pipe-dream models
for Schubert polynomials of the classical groups \cite{SmirnovTutubalina2023},
complementing the divided-difference definitions in types $B$, $C$, and $D$.



\subsection[schubertSkew]{Skew Schubert polynomials}


\begin{polydata}{schubertSkew}
  Name   & Skew Schubert polynomials \\
  Space  & All \\
  Basis  & No \\
  Rating & 1 \\
  Bib    & LenartSottile2003 \\
  Year   & 2003 \\
  Category & Schubert \\
  Contains & schubert | LenartSottile2003 | scope=trivial skewing permutation \\
  PositiveIn & schubert | LenartSottile2003 \\
\end{polydata}

In \cite{LenartSottile2003}, a skew version of Schubert polynomials is
defined.  These polynomials expand positively into Schubert polynomials,
and a formula in the monomial basis is given.
\name{Harry Tamvakis} gives tableau formulas for skew Schubert polynomials
\cite{Tamvakis2023}, including formulas in classical types.


\subsection[schubertDual]{Dual Schubert polynomials}

The dual Schubert polynomials of \name{Alexander Postnikov} and
\name{Richard P. Stanley} are defined by a weighted-chain formula in Bruhat
order \cite{PostnikovStanley2008}.  \name{Zachary Hamaker} gives a
combinatorial proof that this formula is equivalent to the classical Cauchy
identity for Schubert polynomials \cite{Hamaker2023x}.  This also gives a natural
interpretation of the insertion algorithms of \name{Daoji Huang} and
\name{Pavlo Pylyavskyy}.


\subsection[schubertDouble]{Double Schubert polynomials}

\begin{polydata}{schubertDouble}
  Name     & Double Schubert polynomials \\
  Space    & All \\
  Basis    & True \\
  Rating   & 1 \\
  Bib      & LascouxSchutzenberger1982 \\
  Year     & 1982 \\
  Symbol   & $\schubert_\omega(\xvec;\yvec)$ \\
  Keywords & divided-difference  \\
  Category & Schubert \\
  Contains & schurFactorial | ChenLiLouck2002 | scope=Grassmannian permutations and coefficient specialization \\
  SpecializesTo & schubert | LascouxSchutzenberger1982 | map=y=0 \\
\end{polydata}

\begin{polydata}{schubertBackStableDouble}
  Name     & Back stable double Schubert polynomials \\
  Space    & All \\
  Basis    & True \\
  Rating   & 1 \\
  Bib      & LamLeeShimozono2018x \\
  Year     & 2018 \\
  Symbol   & $\overleftarrow{\schubert}_\omega(\xvec;\avec)$ \\
  Keywords & back-stable,double,schubert \\
  Category & Schubert \\
\end{polydata}

\begin{polydata}{schubertEnriched}
  Name     & Enriched Schubert polynomials \\
  Space    & All \\
  Basis    & False \\
  Rating   & 1 \\
  Bib      & AndersonFulton2021x \\
  Year     & 2021 \\
  Keywords & enriched,double,back-stable,type-C,schubert \\
  Category & Schubert \\
  SpecializesTo & schubertDouble | AndersonFulton2021x | map=$c_i=0$ and Anderson--Fulton $\eta_j=-y_j$ \\
  SpecializesTo & schubertBackStableDouble | AndersonFulton2021x | scope=$c$-variables specialized to power series \\
  SpecializesTo & schubertBCD | AndersonFulton2021x | scope=$c$-variables specialized to Schur $Q$-polynomials, type $C$ double case \\
\end{polydata}

The double Schubert polynomials were introduced by
\name{Alain Lascoux} and \name[M.-P. Schützenberger]{Marcel-Paul Schützenberger}
\cite{LascouxSchutzenberger1982}, and by \name{Ian Macdonald} \cite{Macdonald1991}.

They are defined as
\begin{equation*}
 \schubert_\omega(x_1,\dotsc,x_n,y_1,\dotsc,y_n)
 =
 \partial_{\omega^{-1}\omega_0} \prod_{i + j \leq n} (x_i-y_j),
\end{equation*}
where the divided difference operator acts on the $\xvec$-variables.
Setting $y_i=0$ recovers the usual Schubert polynomials.
The double Schubert polynomials generalize the double Schur polynomials, and
the \hyperref[schurFactorial]{factorial Schur polynomials} are obtained from
the Grassmannian double Schur subfamily by specializing the coefficient
alphabet \cite{ChenLiLouck2002}.

\name{David Anderson} and \name{William Fulton} construct enriched type $A$
Schubert polynomials with coefficients in a polynomial ring
\cite{AndersonFulton2021x}.  Specializations recover double Schubert
polynomials, back-stable double Schubert polynomials, and type $C$ double
Schubert polynomials.  \name{David Anderson} also studies Schubert
polynomials through infinite-dimensional flag varieties and degeneracy loci
\cite{Anderson2025}, including positivity in equivariant coproduct and
back-stable/enriched Schubert expansions.

\begin{example}
Anderson and Fulton use the convention
$\schubert^{\mathrm{en}}_w(0,x,\eta)=\schubert_w(x,-\eta)$ for the double
Schubert specialization \cite[Thm.~1.1]{AndersonFulton2021x}.  Thus, to match
the double Schubert convention of this page, we set $\eta_j=-y_j$ after
specializing the enriched variables $c_i$ to zero.  For instance, their
Examples~1 give
\[
  \schubert^{\mathrm{en}}_{s_k}(c,x,\eta)
  =
  c_1+\sum_{i=1}^k (x_i+\eta_i),
  \qquad k>0.
\]
Hence
\[
  \schubert^{\mathrm{en}}_{s_k}(0,x,-y)
  =
  \sum_{i=1}^k (x_i-y_i)
  =
  \schubert_{s_k}(x,y),
\]
the familiar double Schubert polynomial for a simple transposition.
\end{example}

There is a pipe-dream formula, expressing $\schubert_\omega(\xvec,\yvec)$
as a sum over \hyperref[schubertPipeDream]{pipe dreams}:
\begin{equation*}
 \schubert_\omega(\xvec,\yvec) = \sum_{D \in RC(\omega) } \prod_{(i,j) \in D} \left(x_{i} - y_{j} \right).
\end{equation*}
See also the survey by \name{Nantel Bergeron} and \name{Sara Billey} \cite{BergeronBilley1993}.

\name{Allen Knutson} \cite{Knutson2019} provides new proofs of
several combinatorial formulas for the double Schubert polynomials.

\name{Matthew J. Samuel} proves a Molev--Sagan-type formula for products
of double Schubert polynomials under descent hypotheses \cite{Samuel2024}.
The same paper gives a Pieri rule for multiplying a double Schubert polynomial
by a factorial elementary symmetric polynomial, and gives positivity results
after specializing the two equivariant alphabets to be equal.
\name{Yibo Gao} and \name{Rui Xiong} prove Graham positivity results for
triple Schubert calculus \cite{GaoXiong2025x}.  As a corollary, they prove
Kirillov's conjectured positivity for skew divided difference operators
applied to Schubert polynomials.
\name{Serena An}, \name{Katherine Tung}, and \name{Yuchong Zhang} prove that
Postnikov--Stanley polynomials are \hyperref[lorentzianPolynomial]{Lorentzian}
\cite{AnTungZhang2024Lorentzian}.  These polynomials generalize skew dual
Schubert polynomials to arbitrary Weyl groups, and the proof realizes them as
degree polynomials of \hyperref[richardsonVarieties]{Richardson varieties}.


\subsubsection[schubertDoubleGiambelliFormula]{Giambelli formula}

The following Giambelli, or splitting, formula expresses a double Schubert
polynomial in terms of ordinary Schubert polynomials in the two alphabets; see
\cite{Macdonald1991}.  It states that
\begin{equation*}
 \schubert_w(\xvec;\yvec) = \sum_{\substack{ v^{-1}u=w \\ \length(u)+\length(v) = \length(w) }}
 (-1)^{\length(v)}\schubert_u(\xvec) \schubert_v(\yvec).
\end{equation*}
The \hyperref[schubertCauchyFormula]{Cauchy identity} follows by applying this
formula to the longest permutation.



\subsection[schubertAffine]{Affine Schubert polynomials}

\begin{polydata}{schubertAffine}
  Name   & Affine Schubert polynomials \\
  Space  & All \\
  Basis  & No \\
  Rating & 1 \\
  Symbol & $\widetilde{\schubert}_w$ \\
  Bib    & Lam2006Schubert \\
  Keywords & affine, schubert \\
  Year   & 2006 \\
  Category & Schubert \\
  Contains & schurAffine | Lee2015 | scope=$0$-Grassmannian affine permutations \\
  SpecializesTo & stanleySymAffine | Lee2015 | map=$x_i=0$ \\
  SpecializesTo & schubert | Lee2015 | map=$p_j=0$; scope=finite permutations \\
\end{polydata}


In \cite{Lee2015}, an affine extension of Schubert polynomials (indexed by
affine permutations) is defined via divided difference operators after being
introduced by \name{Thomas Lam} in \cite{Lam2006Schubert}.

Fix $n$ and set
\[
  R_n
  =
  \setQ[p_1,\dotsc,p_{n-1}]
  \otimes_{\setQ}
  \setQ[x_1,\dotsc,x_n]/\langle e_1(x),\dotsc,e_n(x)\rangle .
\]
Here the $p_j$ variables come from the affine Grassmannian factor, while the
$x_i$ variables come from the ordinary
\hyperref[flagVarieties]{flag-variety} factor.  Lee defines affine divided
difference operators $\partial_i$, indexed by $i\in \setZ/n\setZ$, on $R_n$.
For $i\neq 0$, these operators fix each $p_m$ and restrict to the usual
divided differences on the $x$-variables.  The affine operator is determined by
\[
  \partial_0(p_m)
  =
  \sum_{j=0}^{m-1} x_1^{m-1-j}x_0^j,
  \qquad
  \partial_i(x_j)=\delta_{ij}-\delta_{i,j+1},
\]
with $x$-indices read modulo $n$.

The \defin{affine Schubert polynomial}\label{affineSchubertPolynomial}
$\widetilde{\schubert}_w$ is the unique
homogeneous element of $R_n$ satisfying
\[
  \partial_i \widetilde{\schubert}_w
  =
  \begin{cases}
    \widetilde{\schubert}_{ws_i}, & \text{if } \length(ws_i)=\length(w)-1,\\
    0, & \text{otherwise,}
  \end{cases}
  \qquad
  \widetilde{\schubert}_{\mathrm{id}}=1
\]
for all $i\in\setZ/n\setZ$ \cite[Thm.~1.2]{Lee2015}.  The projection
$x_i\mapsto 0$ gives the affine Stanley symmetric function, while the projection
$p_j\mapsto 0$ gives the ordinary Schubert polynomial for finite permutations
\cite[Thm.~1.3]{Lee2015}.

\begin{example}
For $n=3$, Lee's first examples are
\[
  \widetilde{\schubert}_{\mathrm{id}}=1,\qquad
  \widetilde{\schubert}_{s_0}=p_1,\qquad
  \widetilde{\schubert}_{s_1}=p_1+x_1,\qquad
  \widetilde{\schubert}_{s_2}=p_1+x_1+x_2 .
\]
Setting $p_1=0$ recovers the ordinary Schubert polynomials
$\schubert_{s_1}=x_1$ and $\schubert_{s_2}=x_1+x_2$.
\end{example}

In \cite[Thm.~6.1]{LamLeeShimozono2019}, the authors prove that the affine Schubert polynomials
expand positively in the monomial basis.


\subsection[schubertInvolution]{Involution Schubert polynomials}

\begin{polydata}{schubertInvolution}
 Name & Involution Schubert polynomials \\
  Space  & All \\
  Basis  & No \\
  Rating & 1 \\
  Symbol & $\ischubert_y(x_1,\dotsc,x_n)$ \\
  Bib    & HamakerMarbergPawlowski2018 \\
  Keywords & divided-difference, schubert \\
  Year   & 2018 \\
  Category & Schubert \\
  StableLimit & stanleySymInvolution | HamakerMarbergPawlowski2018 \\
\end{polydata}

The involution Schubert polynomials were introduced in \cite{WyserYoung2016}.
Additional properties are proved in \cite{HamakerMarbergPawlowski2018},
where the authors introduce the name of the family of polynomials.
The following definition is from that paper.

Let $y \in I_n$, the set of involutions in $\symS_n$, and let $A(y)$ denote the set of
minimal-length permutations $w$ such that $y = w^{-1} \circ w$.
That is, we consider all minimal factorizations of $y$ into simple transpositions
of the form $(s_{i_1}\dotsc s_{i_k})\circ (s_{i_k}\dotsc s_{i_1})$.
Then let
\[
\ischubert_y(x_1,\dotsc,x_n) \coloneqq \sum_{w \in A(y)} \schubert_w(x_1,\dotsc,x_n)
\]
These polynomials can also be produced recursively via divided difference operators.

The stable limit of these yields the
\hyperref[stanleySymInvolution]{involution Stanley symmetric functions}.

A combinatorial model using
\defin{involution pipe dreams}\label{involutionPipeDream} for these polynomials
is given by \name{Zachary Hamaker}, \name{Eric Marberg}, and
\name{Brendan Pawlowski} in \cite{HamakerMarbergPawlowski2019}.

\begin{example}[A small involution Schubert polynomial]
% Related Rust: rust/sym-poly/multipoly/examples/schubert_site_example.rs
For the involution $321 \in \symS_3$, Hamaker--Marberg--Pawlowski give
\[
\hat{R}(321)=\{(s_1,s_2),(s_2,s_1)\},
\qquad
A(321)=\{231,312\}
\]
\cite[Ex.~2.7]{HamakerMarbergPawlowski2018}.  Hence
\[
\ischubert_{321}
  = \schubert_{231}+\schubert_{312}
  = x_1x_2+x_1^2.
\]
The involution pipe-dream model gives a graphical refinement of the same
positive expansion.
\end{example}


\subsection[schubertQuantum]{Quantum Schubert polynomials}


\begin{polydata}{schubertQuantum}
 Name & Quantum Schubert polynomials \\
  Space    & All \\
  Basis    & Yes \\
  Rating   & 3 \\
  Symbol   & $\schubert^q_\omega(\xvec)$ \\
  Bib      & FominGelfandPostnikov1997 \\
  Keywords & schubert \\
  Year     & 1997 \\
  Category & Schubert \\
  SpecializesTo & schubert | FominGelfandPostnikov1997 | map=q=0 \\
\end{polydata}


The quantum Schubert polynomials $\schubert^q_\omega(\xvec)$
are deformations of the Schubert polynomials by a vector $q=(q_1,\dotsc,q_{n-1})$.
These were introduced in \cite{FominGelfandPostnikov1997}.

Recall the \hyperref[schubertAsSumOfElementary]{formula}
that expresses the Schubert polynomials as sums of products of
\hyperref[elementaryE]{elementary symmetric functions}:
\[
\schubert_\omega(\xvec)
=
\sum a_{k_1 \dotsc k_n}
\elementaryE_{k_1}(1)\elementaryE_{k_2}(2) \dotsm \elementaryE_{k_n}(n)
\]
The quantum Schubert polynomials are then defined as
\[
\schubert^q_\omega(\xvec)
=
\sum a_{k_1 \dotsc k_n}
\elementaryE^{q}_{k_1}(1)
\elementaryE^{q}_{k_2}(2)
\dotsm
\elementaryE^{q}_{k_n}(n)
\]
where
\[
\sum_{i=0}^k \elementaryE^{q}_{i}(k)t^i = \det(I+tG_k),
\qquad
G_k=
\begin{pmatrix}
x_1 & q_1 & 0 & \dotsc & 0 \\
-1 & x_2 & q_2 & \dotsc & 0 \\
0 & -1 & x_3 & \ddots & 0 \\
\vdots & \vdots & \ddots & \ddots & q_{k-1} \\
0 & 0 & 0 & -1 & x_k
\end{pmatrix}
\]
Setting $q_i=0$ recovers the classical Schubert polynomials.


There is a bumpless pipe-dream model for the quantum double Schubert polynomials,
see \cite[Thm.~3.3]{LeOuyangTaoRestivoZhang2024x}.
\name{Laura Colmenarejo} and \name{Nicholas Mayers} study the quantum
$k$-Bruhat order that appears in quantum Schubert multiplication
\cite{ColmenarejoMayers2026x}.  They encode maximal chains in intervals by
operators from a free monoid acting on a $q$-extension of $\symS_n$, and
formulate a completeness conjecture for the resulting operator equivalences.


\subsubsection[schubertQuantumMonksRule]{Quantum Monk's rule}

The following refinement of Monk's rule is given in
\cite[Thm.~7.1]{FominGelfandPostnikov1997}.
Let $s_r$ be a simple transposition. Then
\[
\schubert^q_{s_r}(\xvec) \schubert^q_\omega(\xvec) = 
(x_1+x_2+\dotsb + x_r) \schubert^q_\omega(\xvec) 
=
\sum \schubert^q_{\omega t_{ab}}(\xvec) +
\sum q_{cd} \schubert^q_{\omega t_{cd}}(\xvec)
\]
where the first sum is over all transpositions $t_{ab}$ such that $a\leq r \lt b$ and
$\length(\omega t_{ab}) = \length(\omega)+1$,
and the second sum runs over all
transpositions $t_{cd}$ such that $c\leq r \lt d$ and
$\length(\omega t_{cd}) = \length(\omega)-\length(t_{cd}) =
\length(\omega)-2(d-c)+1$.

Here, $q_{cd} = q_c q_{c+1} \dotsm q_{d-1}$.


\subsubsection[schubertQuantumPieriRule]{Quantum Pieri rule}

\name{Sergey Fomin}, \name{Sergey Gelfand}, and \name{Alexander Postnikov}
derive quantum Pieri-type formulas as part of their construction of quantum
Schubert polynomials \cite{FominGelfandPostnikov1997}.
For the Grassmannian subcase, the corresponding quantum multiplication of
Schur polynomials is described by \name{Aaron Bertram},
\name[Ionuț Ciocan-Fontanine]{Ionu{\c{t}} Ciocan-Fontanine}, and
\name{William Fulton} \cite{BertramCiocanFontanineFulton1999}.

\subsubsection[schubertQuantumSottileRule]{Quantum Sottile rule}

\name{Carolina Benedetti Velásquez}, \name{Nantel Bergeron},
\name{Laura Colmenarejo}, \name{Franco Saliola}, and \name{Frank Sottile}
prove a quantum hook Schur multiplication rule in the small quantum
cohomology ring of the flag manifold
\cite[Thm.~3.2]{BenedettiVelasquezBergeronColmenarejoSaliolaSottile2025}.
This is the quantum analogue of the hook/Sottile rule and is the main input
for their quantum Murnaghan--Nakayama rule below.

\subsubsection[schubertQuantumMurnaghanNakayama]{Quantum Murnaghan--Nakayama}

\name{Carolina Benedetti Velásquez}, \name{Nantel Bergeron},
\name{Laura Colmenarejo}, \name{Franco Saliola}, and \name{Frank Sottile}
prove a quantum Murnaghan--Nakayama rule for the flag manifold
\cite{BenedettiVelasquezBergeronColmenarejoSaliolaSottile2025}.  The result
describes multiplication by a tautological class in small quantum cohomology
and is obtained from a hook quantum Schur multiplication formula together
with a detailed analysis of the quantum Bruhat order.




\subsection[schubertQuantumDouble]{Quantum double Schubert polynomials}


\begin{polydata}{schubertQuantumDouble}
 Name & Quantum Double Schubert polynomials \\
  Space    & All \\
  Basis    & Yes \\
  Rating   & 1 \\
  Symbol   & $\schubert^q_\omega(\xvec;\yvec)$ \\
  Bib      & KirillovMaeno2000 \\
  Keywords & schubert \\
  Year     & 2000 \\
  Category & Schubert \\
\end{polydata}


The double quantum Schubert polynomials were introduced by
\name{Anatoly Kirillov} and \name{Toshiaki Maeno}
\cite{KirillovMaeno2000}.


\subsubsection[schubertQuantumDoubleCauchyIdentity]{Cauchy identity}

In \cite{KirillovMaeno2000}, the following Cauchy identity for
quantum double Schubert polynomials is given:
\[
\sum_{\omega \in \symS_n} \schubert^q_\omega(\xvec) \schubert^q_{\omega \omega_0}(\yvec)
= \schubert^q_{\omega_0}(\xvec;\yvec)
\]


\subsection[schubertUniversal]{Universal Schubert polynomials}


\begin{polydata}{schubertUniversal}
  Name & Universal Schubert polynomials \\
  Space    & All \\
  Basis    & Yes \\
  Rating   & 1 \\
  Symbol   & $\schubert_w(c;d)$ \\
  Bib      & Fulton1999 \\
  Keywords & schubert \\
  Year     & 1999 \\
  Category & Schubert \\
  SpecializesTo & schubert | Fulton1999 \\
\end{polydata}


The universal Schubert polynomials were introduced by \name{William Fulton}
\cite{Fulton1999} and
generalize the (double) quantum Schubert polynomials.


The double universal Schubert polynomials are defined via
\begin{equation*}
 \schubert_w(c;d) = \sum_{u,v} (-1)^{\length(v)}\schubert_u(c) \schubert_v(d)
\end{equation*}
as for \hyperref[schubertDoubleGiambelliFormula]{double Schubert polynomials}.


\subsection[schubertTwisted]{Twisted Schubert polynomials}

\begin{polydata}{schubertTwisted}
 Name & Twisted Schubert polynomials \\
  Space    & All \\
  Basis    & Yes \\
  Rating   & 1 \\
  Symbol   & $\tilde{\schubert}_\omega$ \\
  Bib      & Liu2019  \\
  Keywords & schubert \\
  Year     & 2019 \\
  Category & Schubert \\
\end{polydata}

The \defin{twisted Schubert polynomials}\label{twistedSchubertPolynomial}
are defined via $\tilde{\schubert}_{\omega_0}=x_1^{n-1} x_2^{n-2}\dotsm x_{n-1}$,
and the recursion $\tilde{\schubert}_{\omega s_i}= (s_i + \partial_i)\tilde{\schubert}_{\omega}$.

It was proved in \cite{Liu2019} that these are positive in the monomial basis.
This appears to be an early reference that explicitly defines these polynomials.
They are implicitly studied in earlier works, and are related to
the Chern--Schwartz--MacPherson classes of Schubert cells in
\hyperref[flagVarieties]{flag varieties}.
